Wasser siedet, wenn sein Sättigungsdampfdruck gleich dem Druck über der Flüssigkeit ist. Der absolute Druck im Topf ist mit dem äusseren Luftdruck $p_0 = \ncpo$
\begin{align}
 p &= \pTF = \ncpo + \pUe = \pT
\end{align}
Nach der Tabelle für gesättigten Wasserdampf (FoTaBe S.~192) liegt dieser Druck zwischen $p_1 = \psAO$ bei $\vartheta_1 = \thAO$ und $p_2 = \psBO$ bei $\vartheta_2 = \thBO$. Lineare Interpolation ergibt
\begin{align}
 \vartheta &= \thSF \\
   &= \thA + \left(\thB-\thA\right)\cdot\frac{\pTP{3}{4}-\psAO}{\psBO-\psAO} \\
   &= \thS \approx \result{\thSP{0}{3}}
\end{align}
Weil die Dampfdruckkurve nach oben gekrümmt ist, liefert die lineare Interpolation eine etwas zu tiefe Temperatur.
